\chapter{\protect\hyperlink{:}{量子力学}}
\addtocontents{toc}{\protect\hypertarget{:}{}}





\section{\protect\hyperlink{:}{一维定态薛定谔方程}}
\addtocontents{toc}{\protect\hypertarget{:}{}}

量子力学的核心任务是求解薛定谔方程。对于定态问题（势能不显含时间），我们可以将时间部分分离，只需求解\textbf{定态薛定谔方程}——一个关于空间坐标的本征值问题。

\subsection{\protect\hyperlink{:}{定态薛定谔方程的建立}}
\addtocontents{toc}{\protect\hypertarget{:}{}}

含时薛定谔方程为
\begin{equation}
    i\hbar\frac{\partial}{\partial t}\Psi(\vec{r},t) = \hat{H}\Psi(\vec{r},t)
\end{equation}

当哈密顿算符 $\hat{H}$ 不显含时间时（势能 $V(\vec{r})$ 只是位置的函数），我们可以分离变量，令
\begin{equation}
    \Psi(\vec{r},t) = \psi(\vec{r})f(t)
\end{equation}

代入后可以得到
\begin{equation}
    i\hbar\frac{1}{f}\frac{\mathrm{d}f}{\mathrm{d}t} = \frac{1}{\psi}\hat{H}\psi = E
\end{equation}

左边只依赖时间，右边只依赖空间，两边必须等于同一个常数 $E$（能量本征值）。时间部分的解为
\begin{equation}
    f(t) = e^{-iEt/\hbar}
\end{equation}

空间部分满足\textbf{定态薛定谔方程}
\begin{equation}
    \boxed{\hat{H}\psi(\vec{r}) = E\psi(\vec{r})}
\end{equation}

一维情况下为
\begin{equation}
    -\frac{\hbar^2}{2m}\frac{\mathrm{d}^2\psi}{\mathrm{d}x^2} + V(x)\psi = E\psi
\end{equation}

\subsection{\protect\hyperlink{:}{一维无限深方势阱}}
\addtocontents{toc}{\protect\hypertarget{:}{}}

一维无限深方势阱是最简单的量子力学模型，但它揭示了量子力学的许多基本特征。

\begin{definition}[一维无限深方势阱]
    势能函数为
    \begin{equation}
        V(x) = \begin{cases} 0, & 0 < x < a \\ \infty, & x \leq 0 \text{ 或 } x \geq a \end{cases}
    \end{equation}
\end{definition}

由于势阱外的势能为无穷大，粒子不可能出现在阱外，所以边界条件为
\begin{equation}
    \psi(0) = 0, \qquad \psi(a) = 0
\end{equation}

在阱内（$0 < x < a$），$V = 0$，定态方程变为
\begin{equation}
    -\frac{\hbar^2}{2m}\frac{\mathrm{d}^2\psi}{\mathrm{d}x^2} = E\psi
\end{equation}

令 $k^2 = 2mE/\hbar^2$，通解为
\begin{equation}
    \psi(x) = A\sin(kx) + B\cos(kx)
\end{equation}

由边界条件 $\psi(0) = 0$ 得 $B = 0$；由 $\psi(a) = 0$ 得
\begin{equation}
    \sin(ka) = 0 \quad \Longrightarrow \quad ka = n\pi, \quad n = 1, 2, 3, \ldots
\end{equation}

因此能量本征值为
\begin{equation}
    \boxed{E_n = \frac{n^2\pi^2\hbar^2}{2ma^2}, \qquad n = 1, 2, 3, \ldots}
\end{equation}

归一化的本征函数为
\begin{equation}
    \psi_n(x) = \sqrt{\frac{2}{a}}\sin\left(\frac{n\pi x}{a}\right)
\end{equation}

\begin{remark}
    无限深势阱的结果揭示了量子力学的几个重要特征：
    \begin{enumerate}
        \item \textbf{能量量子化}：能量只能取离散的值 $E_n$，这是边界条件的自然结果，不需要人为假设。
        \item \textbf{零点能}：最低能量 $E_1 = \pi^2\hbar^2/(2ma^2) \neq 0$，粒子不可能静止——这是不确定关系的直接体现。
        \item \textbf{波函数的节点}：第 $n$ 个本征态有 $n-1$ 个节点（不包括边界）。
    \end{enumerate}
\end{remark}

\subsection{\protect\hyperlink{:}{一维谐振子}}
\addtocontents{toc}{\protect\hypertarget{:}{}}

一维谐振子是量子力学中最重要的模型之一。它之所以重要，原因有三：
\begin{enumerate}
    \item 谐振子势 $V(x) = \frac{1}{2}m\omega^2 x^2$ 是任何稳定平衡位置附近势能的通用近似（泰勒展开到二阶）。
    \item 它是可以精确求解的少数量子力学问题之一。
    \item 它在量子场论中扮演核心角色——场的量子化本质上就是无穷多个谐振子的量子化。
\end{enumerate}

\subsubsection{\protect\hyperlink{:}{从经典到量子}}
\addtocontents{toc}{\protect\hypertarget{:}{}}

经典力学中，一个弹簧振子的运动方程为 $m\ddot{x} = -kx$，解为 $x(t) = A\cos(\omega t + \varphi)$，其中 $\omega = \sqrt{k/m}$。能量为
\begin{equation}
    E = \frac{1}{2}kA^2 = \frac{1}{2}m\omega^2 A^2
\end{equation}

经典力学允许能量取任意非负值。量子力学会给出什么不同的结果呢？

\subsubsection{\protect\hyperlink{:}{定态薛定谔方程的建立}}
\addtocontents{toc}{\protect\hypertarget{:}{}}%

\begin{definition}[一维谐振子]
    势能函数为
    \begin{equation}
        V(x) = \frac{1}{2}m\omega^2 x^2
    \end{equation}
    其中 $\omega = \sqrt{k/m}$ 为经典角频率。
\end{definition}

一维谐振子的哈密顿量为
\begin{equation}
    \hat{H} = \frac{\hat{p}^2}{2m} + \frac{1}{2}m\omega^2\hat{x}^2 = -\frac{\hbar^2}{2m}\frac{\mathrm{d}^2}{\mathrm{d}x^2} + \frac{1}{2}m\omega^2 x^2
\end{equation}

定态薛定谔方程为
\begin{equation}
    -\frac{\hbar^2}{2m}\frac{\mathrm{d}^2\psi}{\mathrm{d}x^2} + \frac{1}{2}m\omega^2 x^2\psi = E\psi
\end{equation}

这是一个变系数二阶常微分方程。下面介绍两种求解方法。

\subsubsection{\protect\hyperlink{:}{方法一：幂级数法（解析法）}}
\addtocontents{toc}{\protect\hypertarget{:}{}}

引入无量纲变量
\begin{equation}
    \xi = \sqrt{\frac{m\omega}{\hbar}}\,x, \qquad \lambda = \frac{2E}{\hbar\omega}
\end{equation}

方程化为
\begin{equation}
    \frac{\mathrm{d}^2\psi}{\mathrm{d}\xi^2} + (\lambda - \xi^2)\psi = 0
\end{equation}

\paragraph*{第一步：分析渐近行为}

当 $|\xi| \to \infty$ 时，$\lambda$ 可以忽略，方程近似为 $\psi'' \approx \xi^2\psi$，其渐近解为 $\psi \sim e^{\pm\xi^2/2}$。由于波函数必须可归一化，取负号
\begin{equation}
    \psi(\xi) \sim e^{-\xi^2/2}, \qquad |\xi| \to \infty
\end{equation}

\paragraph*{第二步：提取渐近因子}

令
\begin{equation}
    \psi(\xi) = h(\xi)e^{-\xi^2/2}
\end{equation}

代入原方程，得到 $h(\xi)$ 满足的方程
\begin{equation}
    h'' - 2\xi h' + (\lambda - 1)h = 0
\end{equation}

这就是\textbf{厄米方程}。

\paragraph*{第三步：幂级数展开}

设 $h(\xi)$ 可以展开为幂级数
\begin{equation}
    h(\xi) = \sum_{k=0}^{\infty} a_k \xi^k
\end{equation}

代入厄米方程，比较 $\xi^k$ 的系数，得到递推关系
\begin{equation}
    a_{k+2} = \frac{2k + 1 - \lambda}{(k+1)(k+2)}a_k
\end{equation}

\paragraph*{第四步：量子化条件}

当 $k \to \infty$ 时，$a_{k+2}/a_k \approx 2/k$，这与 $e^{\xi^2}$ 的展开系数相同。如果级数不截断，$\psi \sim h(\xi)e^{-\xi^2/2} \sim e^{+\xi^2/2}$，在无穷远处发散，不可归一化。

因此，级数必须在某一项截断。设最高次项为 $k = n$，则要求
\begin{equation}
    2n + 1 - \lambda = 0 \quad \Longrightarrow \quad \lambda = 2n + 1
\end{equation}

代回 $\lambda = 2E/(\hbar\omega)$，得到\textbf{能量量子化条件}
\begin{equation}
    \boxed{E_n = \left(n + \frac{1}{2}\right)\hbar\omega, \qquad n = 0, 1, 2, \ldots}
\end{equation}

此时 $h(\xi)$ 是一个 $n$ 次多项式——\textbf{厄米多项式} $H_n(\xi)$。

\subsubsection{\protect\hyperlink{:}{方法二：代数法（升降算符）}}
\addtocontents{toc}{\protect\hypertarget{:}{}}

代数法不需要直接求解微分方程，而是利用算符的代数性质来得到能谱。这种方法更加优雅，也更能揭示问题的深层结构。

\paragraph*{第一步：定义升降算符}

将哈密顿量写成
\begin{equation}
    \hat{H} = \frac{\hat{p}^2}{2m} + \frac{1}{2}m\omega^2\hat{x}^2
\end{equation}

我们希望将它因式分解。注意到 $\hat{x}$ 和 $\hat{p}$ 满足对易关系 $[\hat{x}, \hat{p}] = i\hbar$，定义

\begin{definition}[升降算符]
    \begin{equation}
        \hat{a} = \sqrt{\frac{m\omega}{2\hbar}}\left(\hat{x} + \frac{i}{m\omega}\hat{p}\right), \qquad \hat{a}^\dagger = \sqrt{\frac{m\omega}{2\hbar}}\left(\hat{x} - \frac{i}{m\omega}\hat{p}\right)
    \end{equation}
    其中 $\hat{a}$ 称为\textbf{湮灭算符}（或降算符），$\hat{a}^\dagger$ 称为\textbf{产生算符}（或升算符）。
\end{definition}

\paragraph*{第二步：推导对易关系}

利用 $[\hat{x}, \hat{p}] = i\hbar$，可以计算

\begin{theorem}[升降算符的对易关系]
    \begin{equation}
        [\hat{a}, \hat{a}^\dagger] = 1
    \end{equation}
\end{theorem}

\begin{proof}
    \begin{equation}
        [\hat{a}, \hat{a}^\dagger] = \frac{m\omega}{2\hbar}\left[\hat{x} + \frac{i}{m\omega}\hat{p}, \hat{x} - \frac{i}{m\omega}\hat{p}\right] = \frac{m\omega}{2\hbar}\cdot\frac{-2i}{m\omega}[\hat{x}, \hat{p}] = \frac{-2i}{2\hbar}\cdot i\hbar = 1
    \end{equation}
\end{proof}

\paragraph*{第三步：用升降算符表示哈密顿量}

计算 $\hat{a}^\dagger\hat{a}$
\begin{equation}
    \hat{a}^\dagger\hat{a} = \frac{m\omega}{2\hbar}\left(\hat{x} - \frac{i}{m\omega}\hat{p}\right)\left(\hat{x} + \frac{i}{m\omega}\hat{p}\right) = \frac{m\omega}{2\hbar}\hat{x}^2 + \frac{1}{2m\hbar\omega}\hat{p}^2 - \frac{i}{2\hbar}[\hat{x}, \hat{p}]
\end{equation}

整理后得到

\begin{theorem}[哈密顿量的因式分解]
    定义粒子数算符 $\hat{N} = \hat{a}^\dagger\hat{a}$，则哈密顿量可以写为
    \begin{equation}
        \hat{H} = \left(\hat{a}^\dagger\hat{a} + \frac{1}{2}\right)\hbar\omega = \left(\hat{N} + \frac{1}{2}\right)\hbar\omega
    \end{equation}
\end{theorem}

\paragraph*{第四步：确定能谱}

设 $\hat{N}$ 的本征值为 $n$，本征态为 $\ket{n}$，即 $\hat{N}\ket{n} = n\ket{n}$。

首先证明 $n \geq 0$。由于
\begin{equation}
    n = \bra{n}\hat{a}^\dagger\hat{a}\ket{n} = \|\hat{a}\ket{n}\|^2 \geq 0
\end{equation}

等号成立当且仅当 $\hat{a}\ket{n} = 0$。

接下来证明 $\hat{a}\ket{n}$ 是 $\hat{N}$ 的本征态，本征值为 $n-1$
\begin{equation}
    \hat{N}(\hat{a}\ket{n}) = \hat{a}^\dagger\hat{a}\hat{a}\ket{n} = (\hat{a}\hat{a}^\dagger - 1)\hat{a}\ket{n} = \hat{a}(\hat{a}^\dagger\hat{a} - 1)\ket{n} = (n-1)(\hat{a}\ket{n})
\end{equation}

同理，$\hat{a}^\dagger\ket{n}$ 是 $\hat{N}$ 的本征态，本征值为 $n+1$
\begin{equation}
    \hat{N}(\hat{a}^\dagger\ket{n}) = \hat{a}^\dagger\hat{a}\hat{a}^\dagger\ket{n} = \hat{a}^\dagger(\hat{a}^\dagger\hat{a} + 1)\ket{n} = (n+1)(\hat{a}^\dagger\ket{n})
\end{equation}

这意味着 $\hat{a}$ 将本征值降低 1，$\hat{a}^\dagger$ 将本征值升高 1。不断应用 $\hat{a}$ 会得到越来越小的本征值，但由于 $n \geq 0$，这个过程必须终止。终止条件为
\begin{equation}
    \hat{a}\ket{0} = 0
\end{equation}

此时 $n = 0$。从基态出发，不断应用 $\hat{a}^\dagger$，得到所有本征值
\begin{equation}
    n = 0, 1, 2, 3, \ldots
\end{equation}

因此能量本征值为
\begin{equation}
    \boxed{E_n = \left(n + \frac{1}{2}\right)\hbar\omega, \qquad n = 0, 1, 2, \ldots}
\end{equation}

\paragraph*{第五步：升降算符的作用}

\begin{theorem}[升降算符对本征态的作用]
    归一化后，升降算符对本征态的作用为
    \begin{align}
        \hat{a}^\dagger\ket{n} &= \sqrt{n+1}\ket{n+1} \\
        \hat{a}\ket{n} &= \sqrt{n}\ket{n-1}
    \end{align}
\end{theorem}

这就是为什么它们被称为"产生"和"湮灭"算符——产生算符产生一个量子，湮灭算符消灭一个量子。

\subsubsection{\protect\hyperlink{:}{能量本征值与零点能}}
\addtocontents{toc}{\protect\hypertarget{:}{}}

\begin{theorem}[一维谐振子的能谱]
    一维谐振子的能量本征值为
    \begin{equation}
        E_n = \left(n + \frac{1}{2}\right)\hbar\omega, \qquad n = 0, 1, 2, \ldots
    \end{equation}
    能级是等间距的，间距为 $\hbar\omega$。
\end{theorem}

\begin{remark}
    \begin{enumerate}
        \item \textbf{零点能}：最低能量为 $E_0 = \frac{1}{2}\hbar\omega \neq 0$，粒子在基态时仍然具有能量。它的物理根源是不确定关系：如果粒子完全静止（$p = 0$），则 $\Delta p = 0$，由 $\Delta x \Delta p \geq \hbar/2$ 得 $\Delta x \to \infty$，粒子将不受束缚。
        \item \textbf{能量量子化}是边界条件（波函数可归一化）的自然结果，不需要像旧量子论那样人为假设。
        \item \textbf{经典类比}：经典谐振子的平均能量为 $\langle E \rangle = \frac{1}{2}kA^2$，可以取任意非负值。量子谐振子的基态能量 $\frac{1}{2}\hbar\omega$ 对应于经典振子振幅为零时的能量——但量子力学不允许振幅精确为零。
    \end{enumerate}
\end{remark}

\subsubsection{\protect\hyperlink{:}{波函数与概率分布}}
\addtocontents{toc}{\protect\hypertarget{:}{}}

\begin{definition}[基态波函数]
    由 $\hat{a}\ket{0} = 0$，在坐标表象中为
    \begin{equation}
        \left(\hat{x} + \frac{i}{m\omega}\hat{p}\right)\psi_0(x) = 0 \quad \Longrightarrow \quad x\psi_0 + \frac{\hbar}{m\omega}\frac{\mathrm{d}\psi_0}{\mathrm{d}x} = 0
    \end{equation}
    解这个一阶微分方程并归一化后
    \begin{equation}
        \psi_0(x) = \left(\frac{m\omega}{\pi\hbar}\right)^{1/4} e^{-m\omega x^2/(2\hbar)}
    \end{equation}
    这是一个高斯函数，其概率密度为
    \begin{equation}
        |\psi_0(x)|^2 = \sqrt{\frac{m\omega}{\pi\hbar}}\,e^{-m\omega x^2/\hbar}
    \end{equation}
\end{definition}

\paragraph*{激发态波函数} 利用产生算符递推
\begin{equation}
    \ket{n} = \frac{(\hat{a}^\dagger)^n}{\sqrt{n!}}\ket{0}
\end{equation}

在坐标表象中
\begin{equation}
    \psi_n(x) = \left(\frac{m\omega}{\pi\hbar}\right)^{1/4} \frac{1}{\sqrt{2^n n!}} H_n(\xi) e^{-\xi^2/2}
\end{equation}

其中 $\xi = \sqrt{m\omega/\hbar}\,x$，$H_n(\xi)$ 是第 $n$ 阶厄米多项式。

前几个厄米多项式为
\begin{align}
    H_0(\xi) &= 1 \\
    H_1(\xi) &= 2\xi \\
    H_2(\xi) &= 4\xi^2 - 2 \\
    H_3(\xi) &= 8\xi^3 - 12\xi
\end{align}

厄米多项式的递推关系为
\begin{equation}
    H_{n+1}(\xi) = 2\xi H_n(\xi) - 2n H_{n-1}(\xi)
\end{equation}

\paragraph*{波函数的节点} 第 $n$ 个激发态 $\psi_n(x)$ 有 $n$ 个节点（波函数为零的点）。这与无限深势阱的情况类似，是波函数正交性的几何体现。

\paragraph*{宇称} 谐振子势 $V(x) = \frac{1}{2}m\omega^2 x^2$ 关于 $x = 0$ 对称，因此本征函数具有确定的宇称
\begin{equation}
    \psi_n(-x) = (-1)^n \psi_n(x)
\end{equation}

偶数 $n$ 对应偶宇称，奇数 $n$ 对应奇宇称。

\subsubsection{\protect\hyperlink{:}{经典极限与对应原理}}
\addtocontents{toc}{\protect\hypertarget{:}{}}

当量子数 $n$ 很大时，量子谐振子的行为应该趋近于经典谐振子——这就是\textbf{对应原理}的体现。

对于高激发态 $n \gg 1$：
\begin{enumerate}
    \item 能级间距 $\hbar\omega$ 相对于总能量 $E_n \approx n\hbar\omega$ 来说可以忽略，能量近似连续。
    \item 概率密度 $|\psi_n(x)|^2$ 在经典允许区域 $|x| \leq A = \sqrt{2E/(m\omega^2)}$ 内近似均匀分布，在经典禁区边缘有明显的峰值（\textbf{量子隧穿}的体现）。
\end{enumerate}

这种从量子到经典的过渡，是理解量子力学与经典力学关系的重要窗口。

\subsection{\protect\hyperlink{:}{一维方势垒与量子隧穿}}
\addtocontents{toc}{\protect\hypertarget{:}{}}

量子隧穿是量子力学最奇特的效应之一：粒子可以穿越经典力学不允许的势垒区域。

\begin{definition}[一维方势垒]
    势能函数为
    \begin{equation}
        V(x) = \begin{cases} V_0, & 0 < x < a \\ 0, & \text{其他} \end{cases}
    \end{equation}
\end{definition}

设能量 $E < V_0$ 的粒子从左侧入射。经典力学预言粒子一定会被反射，但量子力学给出不同的结果。

在三个区域中，波函数分别为：
\begin{enumerate}
    \item 区域 I（$x < 0$）：$\psi_1 = Ae^{ikx} + Be^{-ikx}$，$k = \sqrt{2mE}/\hbar$
    \item 区域 II（$0 < x < a$）：$\psi_2 = Ce^{-\kappa x} + De^{\kappa x}$，$\kappa = \sqrt{2m(V_0-E)}/\hbar$
    \item 区域 III（$x > a$）：$\psi_3 = Fe^{ikx}$（只有透射波）
\end{enumerate}

利用波函数及其导数在边界处连续的条件，可以求出透射系数
\begin{equation}
    T = \frac{|F|^2}{|A|^2} = \frac{1}{1 + \frac{V_0^2\sinh^2(\kappa a)}{4E(V_0-E)}}
\end{equation}

当 $\kappa a \gg 1$ 时（宽势垒或高势垒），近似有
\begin{equation}
    T \approx \frac{16E(V_0-E)}{V_0^2}e^{-2\kappa a}
\end{equation}

透射系数随势垒宽度 $a$ 和高度 $V_0$ 指数衰减，但\textbf{不为零}——这就是量子隧穿效应。

\begin{remark}
    量子隧穿在物理中有重要应用：$\alpha$ 衰变、扫描隧道显微镜（STM）、半导体中的隧道二极管等。
\end{remark}

\subsection{\protect\hyperlink{:}{束缚态与散射态}}
\addtocontents{toc}{\protect\hypertarget{:}{}}

根据能量的取值，一维定态问题可以分为两类：

\begin{enumerate}
    \item \textbf{束缚态}：粒子被限制在一个有限区域内运动。能量取离散值（量子化），波函数在无穷远处趋于零。如无限深势阱和谐振子。
    \item \textbf{散射态}：粒子可以从无穷远处来，又回到无穷远处。能量可以取任意正值（连续谱），波函数在无穷远处为平面波。如方势垒散射。
\end{enumerate}

束缚态和散射态的区别是量子力学中一个重要的概念，它反映了边界条件对物理系统的深刻影响。



\section{\protect\hyperlink{:}{算符}}
\addtocontents{toc}{\protect\hypertarget{:}{}}

在量子力学中，可观测量（力学量）用算符表示。算符是作用在希尔伯特空间上的线性变换。本节系统介绍厄米算符的性质、重要定理以及常用结论。

\subsection{\protect\hyperlink{:}{基本概念}}
\addtocontents{toc}{\protect\hypertarget{:}{}}

\subsubsection{\protect\hyperlink{:}{算符的定义与运算}}
\addtocontents{toc}{\protect\hypertarget{:}{}}

\begin{definition}[线性算符]
    算符 $\hat{A}$ 是将一个态矢量映射为另一个态矢量的线性变换，即对任意态矢量 $\ket{\psi}$、$\ket{\phi}$ 和复数 $c_1$、$c_2$
    \begin{equation}
        \hat{A}(c_1\ket{\psi} + c_2\ket{\phi}) = c_1\hat{A}\ket{\psi} + c_2\hat{A}\ket{\phi}
    \end{equation}
\end{definition}

算符的基本运算包括：
\begin{enumerate}
    \item \textbf{乘积}：$(\hat{A}\hat{B})\ket{\psi} = \hat{A}(\hat{B}\ket{\psi})$，一般 $\hat{A}\hat{B} \neq \hat{B}\hat{A}$。
    \item \textbf{对易子}：$[\hat{A}, \hat{B}] = \hat{A}\hat{B} - \hat{B}\hat{A}$。若 $[\hat{A}, \hat{B}] = 0$，则称 $\hat{A}$ 和 $\hat{B}$ 对易。
    \item \textbf{反对易子}：$\{\hat{A}, \hat{B}\} = \hat{A}\hat{B} + \hat{B}\hat{A}$。
\end{enumerate}

\subsubsection{\protect\hyperlink{:}{厄米算符}}
\addtocontents{toc}{\protect\hypertarget{:}{}}

\begin{definition}[伴随算符（共轭算符）]
    算符 $\hat{A}$ 的伴随算符 $\hat{A}^\dagger$ 定义为：对任意态矢量 $\ket{\psi}$、$\ket{\phi}$
    \begin{equation}
        \bra{\psi}\hat{A}^\dagger\ket{\phi} = \bra{\phi}\hat{A}\ket{\psi}^*
    \end{equation}
\end{definition}

\begin{definition}[厄米算符（自伴算符）]
    若算符 $\hat{A}$ 满足 $\hat{A} = \hat{A}^\dagger$，则称 $\hat{A}$ 为厄米算符。等价地，对任意态矢量 $\ket{\psi}$、$\ket{\phi}$
    \begin{equation}
        \bra{\psi}\hat{A}\ket{\phi} = \bra{\phi}\hat{A}\ket{\psi}^*
    \end{equation}
\end{definition}

\begin{remark}
    量子力学中，所有可观测量都对应厄米算符。这是因为厄米算符的本征值为实数，而物理测量的结果必须是实数。
\end{remark}

\subsection{\protect\hyperlink{:}{厄米算符的性质}}
\addtocontents{toc}{\protect\hypertarget{:}{}}

厄米算符具有一系列重要的数学性质，这些性质保证了量子力学的自洽性。

\subsubsection{\protect\hyperlink{:}{本征值为实数}}
\addtocontents{toc}{\protect\hypertarget{:}{}}

\begin{theorem}[厄米算符的本征值为实数]
    设 $\hat{A}$ 为厄米算符，$\ket{\psi}$ 为其本征态，本征值为 $a$，即 $\hat{A}\ket{\psi} = a\ket{\psi}$，则 $a$ 必为实数。
\end{theorem}

\begin{proof}
    由本征方程 $\hat{A}\ket{\psi} = a\ket{\psi}$，左乘 $\bra{\psi}$
    \begin{equation}
        \bra{\psi}\hat{A}\ket{\psi} = a\bra{\psi}\ket{\psi}
    \end{equation}
    取复共轭
    \begin{equation}
        \bra{\psi}\hat{A}\ket{\psi}^* = a^*\bra{\psi}\ket{\psi}^*
    \end{equation}
    由于 $\hat{A} = \hat{A}^\dagger$，有 $\bra{\psi}\hat{A}\ket{\psi}^* = \bra{\psi}\hat{A}^\dagger\ket{\psi} = \bra{\psi}\hat{A}\ket{\psi}$，因此
    \begin{equation}
        a\bra{\psi}\ket{\psi} = a^*\bra{\psi}\ket{\psi}
    \end{equation}
    由于 $\bra{\psi}\ket{\psi} > 0$（非零态矢量的模方为正），故 $a = a^*$，即 $a$ 为实数。
\end{proof}

\subsubsection{\protect\hyperlink{:}{不同本征值对应的本征态正交}}
\addtocontents{toc}{\protect\hypertarget{:}{}}

\begin{theorem}[厄米算符本征态的正交性]
    设 $\hat{A}$ 为厄米算符，$\ket{\psi_1}$ 和 $\ket{\psi_2}$ 分别是本征值 $a_1$ 和 $a_2$ 对应的本征态，若 $a_1 \neq a_2$，则 $\ket{\psi_1}$ 与 $\ket{\psi_2}$ 正交，即 $\braket{\psi_1}{\psi_2} = 0$。
\end{theorem}

\begin{proof}
    由本征方程
    \begin{equation}
        \hat{A}\ket{\psi_1} = a_1\ket{\psi_1}, \qquad \hat{A}\ket{\psi_2} = a_2\ket{\psi_2}
    \end{equation}
    计算 $\bra{\psi_1}\hat{A}\ket{\psi_2}$ 两种方式：
    \begin{enumerate}
        \item 将 $\hat{A}$ 作用于右边：$\bra{\psi_1}\hat{A}\ket{\psi_2} = a_2\braket{\psi_1}{\psi_2}$
        \item 利用厄米性，将 $\hat{A}$ 作用于左边：$\bra{\psi_1}\hat{A}\ket{\psi_2} = (\bra{\psi_2}\hat{A}\ket{\psi_1})^* = a_1^*\braket{\psi_1}{\psi_2}^* = a_1\braket{\psi_2}{\psi_1}$
    \end{enumerate}
    两式相减
    \begin{equation}
        (a_2 - a_1)\braket{\psi_1}{\psi_2} = 0
    \end{equation}
    由于 $a_1 \neq a_2$，故 $\braket{\psi_1}{\psi_2} = 0$。
\end{proof}

若本征值有简并（同一本征值对应多个线性无关的本征态），可以通过施密特正交化方法构造正交归一的本征态系。因此，厄米算符的本征态总可以选为正交归一的
\begin{equation}
    \braket{\psi_i}{\psi_j} = \delta_{ij}
\end{equation}

\subsubsection{\protect\hyperlink{:}{本征态的完备性}}
\addtocontents{toc}{\protect\hypertarget{:}{}}

\begin{theorem}[厄米算符本征态的完备性]
    厄米算符的本征态构成希尔伯特空间的一组完备基。即任意态矢量 $\ket{\Psi}$ 都可以展开为
    \begin{equation}
        \ket{\Psi} = \sum_n c_n\ket{\psi_n}, \qquad c_n = \braket{\psi_n}{\Psi}
    \end{equation}
    等价地，完备性关系为
    \begin{equation}
        \sum_n \ket{\psi_n}\bra{\psi_n} = \hat{1}
    \end{equation}
    其中 $\hat{1}$ 是单位算符。
\end{theorem}

\begin{remark}
    完备性是一个物理上非常强的要求——它保证了任意量子态都可以用某组本征态来展开。严格来说，完备性需要证明，但在物理应用中通常作为基本假设接受。
\end{remark}

\subsubsection{\protect\hyperlink{:}{对易算符的共同本征态}}
\addtocontents{toc}{\protect\hypertarget{:}{}}

\begin{theorem}[对易厄米算符有共同本征态系]
    若两个厄米算符 $\hat{A}$ 和 $\hat{B}$ 对易（$[\hat{A}, \hat{B}] = 0$），则它们存在共同的本征态系。
\end{theorem}

\begin{proof}
    设 $\ket{a}$ 是 $\hat{A}$ 的本征态，本征值为 $a$。考虑 $\hat{B}\ket{a}$
    \begin{equation}
        \hat{A}(\hat{B}\ket{a}) = \hat{B}\hat{A}\ket{a} = a(\hat{B}\ket{a})
    \end{equation}
    因此 $\hat{B}\ket{a}$ 也是 $\hat{A}$ 的本征态，本征值同为 $a$。
    \begin{enumerate}
        \item 若 $a$ 是非简并的，则 $\hat{B}\ket{a}$ 必与 $\ket{a}$ 成比例，即 $\hat{B}\ket{a} = b\ket{a}$，$\ket{a}$ 也是 $\hat{B}$ 的本征态。
        \item 若 $a$ 是简并的，可以在简并子空间内对角化 $\hat{B}$，找到的本征态同时是 $\hat{A}$ 和 $\hat{B}$ 的本征态。
    \end{enumerate}
\end{proof}

\begin{definition}[完全对易观测量集（CSCO）]
    若一组相互对易的厄米算符 $\{\hat{A}, \hat{B}, \hat{C}, \ldots\}$ 的共同本征态可以完全确定系统的量子态（即不存在简并），则称这组算符为\textbf{完全对易观测量集}（CSCO）。
\end{definition}

例如，氢原子中 $\{\hat{H}, \hat{L}^2, \hat{L}_z\}$ 构成 CSCO，共同本征态为 $\ket{n,l,m}$。

\subsection{\protect\hyperlink{:}{重要定理}}
\addtocontents{toc}{\protect\hypertarget{:}{}}

\subsubsection{\protect\hyperlink{:}{不确定关系}}
\addtocontents{toc}{\protect\hypertarget{:}{}}

不确定关系是量子力学最基本的定理之一，它揭示了量子力学中测量精度的根本限制。

\begin{definition}[标准差]
    对于力学量 $\hat{A}$，在态 $\ket{\Psi}$ 中的标准差定义为
    \begin{equation}
        \sigma_A = \sqrt{\langle \hat{A}^2 \rangle - \langle \hat{A} \rangle^2}
    \end{equation}
    其中 $\langle \hat{A} \rangle = \bra{\Psi}\hat{A}\ket{\Psi}$ 是 $\hat{A}$ 的期望值。
\end{definition}

\begin{theorem}[不确定关系（Robertson不等式）]
    对于任意两个厄米算符 $\hat{A}$ 和 $\hat{B}$，在任意态 $\ket{\Psi}$ 中
    \begin{equation}
        \sigma_A \sigma_B \geqslant \frac{1}{2}\left|\langle [\hat{A}, \hat{B}] \rangle\right|
    \end{equation}
\end{theorem}

\begin{proof}
    定义两个新算符
    \begin{equation}
        \delta\hat{A} = \hat{A} - \langle \hat{A} \rangle, \qquad \delta\hat{B} = \hat{B} - \langle \hat{B} \rangle
    \end{equation}
    显然 $\sigma_A^2 = \langle (\delta\hat{A})^2 \rangle$，$\sigma_B^2 = \langle (\delta\hat{B})^2 \rangle$，且 $[\delta\hat{A}, \delta\hat{B}] = [\hat{A}, \hat{B}]$。

    考虑积分
    \begin{equation}
        I(\lambda) = \int_{-\infty}^{+\infty} |\delta\hat{A}\psi + i\lambda\delta\hat{B}\psi|^2 \mathrm{d}x \geqslant 0
    \end{equation}
    展开
    \begin{equation}
        I(\lambda) = \langle (\delta\hat{A})^2 \rangle + \lambda^2\langle (\delta\hat{B})^2 \rangle + i\lambda\langle [\delta\hat{A}, \delta\hat{B}] \rangle \geqslant 0
    \end{equation}
    即
    \begin{equation}
        \sigma_A^2 + \lambda^2\sigma_B^2 + i\lambda\langle [\hat{A}, \hat{B}] \rangle \geqslant 0
    \end{equation}
    由于 $[\hat{A}, \hat{B}]$ 是反对称的（$[\hat{A}, \hat{B}]^\dagger = -[\hat{A}, \hat{B}]$），$\langle [\hat{A}, \hat{B}] \rangle$ 是纯虚数，令 $\langle [\hat{A}, \hat{B}] \rangle = i\alpha$（$\alpha$ 为实数），则
    \begin{equation}
        \sigma_B^2\lambda^2 - \alpha\lambda + \sigma_A^2 \geqslant 0
    \end{equation}
    这是关于 $\lambda$ 的二次不等式，对任意 $\lambda$ 成立的条件是判别式 $\leqslant 0$
    \begin{equation}
        \alpha^2 - 4\sigma_A^2\sigma_B^2 \leqslant 0 \quad \Longrightarrow \quad \sigma_A^2\sigma_B^2 \geqslant \frac{\alpha^2}{4} = \frac{1}{4}|\langle [\hat{A}, \hat{B}] \rangle|^2
    \end{equation}
    开方即得
    \begin{equation}
        \sigma_A\sigma_B \geqslant \frac{1}{2}|\langle [\hat{A}, \hat{B}] \rangle|
    \end{equation}
\end{proof}

\begin{remark}
    最常用的不确定关系是位置-动量不确定关系。由 $[\hat{x}, \hat{p}] = i\hbar$，得
    \begin{equation}
        \sigma_x\sigma_p \geqslant \frac{\hbar}{2}
    \end{equation}
    等号在高斯波包（谐振子基态）时取到。
\end{remark}

\subsubsection{\protect\hyperlink{:}{埃伦费斯特定理}}
\addtocontents{toc}{\protect\hypertarget{:}{}}

埃伦费斯特定理描述了力学量期望值随时间的演化，它是量子力学与经典力学对应关系的桥梁。

\begin{theorem}[埃伦费斯特定理]
    对于任意力学量 $\hat{F}$（可能显含时间），其期望值随时间的变化率为
    \begin{equation}
        \frac{\mathrm{d}}{\mathrm{d}t}\langle \hat{F} \rangle = \frac{\partial\langle \hat{F} \rangle}{\partial t} + \frac{1}{i\hbar}\langle [\hat{F}, \hat{H}] \rangle
    \end{equation}
\end{theorem}

\begin{proof}
    由定义 $\langle \hat{F} \rangle = \bra{\Psi(t)}\hat{F}\ket{\Psi(t)}$，对时间求导
    \begin{equation}
        \frac{\mathrm{d}}{\mathrm{d}t}\langle \hat{F} \rangle = \left(\frac{\partial\bra{\Psi}}{\partial t}\right)\hat{F}\ket{\Psi} + \bra{\Psi}\frac{\partial\hat{F}}{\partial t}\ket{\Psi} + \bra{\Psi}\hat{F}\left(\frac{\partial\ket{\Psi}}{\partial t}\right)
    \end{equation}
    由薛定谔方程 $i\hbar\frac{\partial\ket{\Psi}}{\partial t} = \hat{H}\ket{\Psi}$，得
    \begin{equation}
        \frac{\partial\ket{\Psi}}{\partial t} = \frac{1}{i\hbar}\hat{H}\ket{\Psi}, \qquad \frac{\partial\bra{\Psi}}{\partial t} = -\frac{1}{i\hbar}\bra{\Psi}\hat{H}
    \end{equation}
    代入
    \begin{align}
        \frac{\mathrm{d}}{\mathrm{d}t}\langle \hat{F} \rangle &= -\frac{1}{i\hbar}\bra{\Psi}\hat{H}\hat{F}\ket{\Psi} + \frac{\partial\langle \hat{F} \rangle}{\partial t} + \frac{1}{i\hbar}\bra{\Psi}\hat{F}\hat{H}\ket{\Psi} \nonumber \\
        &= \frac{\partial\langle \hat{F} \rangle}{\partial t} + \frac{1}{i\hbar}\bra{\Psi}[\hat{F}, \hat{H}]\ket{\Psi} \nonumber \\
        &= \frac{\partial\langle \hat{F} \rangle}{\partial t} + \frac{1}{i\hbar}\langle [\hat{F}, \hat{H}] \rangle
    \end{align}
\end{proof}

\begin{remark}
    埃伦费斯特定理给出了\textbf{守恒量}的判据：若 $\hat{F}$ 不显含时间（$\frac{\partial\hat{F}}{\partial t} = 0$）且与哈密顿量对易（$[\hat{F}, \hat{H}] = 0$），则 $\frac{\mathrm{d}}{\mathrm{d}t}\langle \hat{F} \rangle = 0$，即 $\hat{F}$ 的期望值不随时间变化。例如，在中心力场中，$\hat{L}^2$ 和 $\hat{L}_z$ 都是守恒量。
\end{remark}

\begin{example}[经典极限]
    取 $\hat{F} = \hat{x}$ 和 $\hat{F} = \hat{p}$，由埃伦费斯特定理
    \begin{align}
        \frac{\mathrm{d}}{\mathrm{d}t}\langle \hat{x} \rangle &= \frac{1}{i\hbar}\langle [\hat{x}, \hat{H}] \rangle = \frac{\langle \hat{p} \rangle}{m} \\
        \frac{\mathrm{d}}{\mathrm{d}t}\langle \hat{p} \rangle &= \frac{1}{i\hbar}\langle [\hat{p}, \hat{H}] \rangle = -\left\langle \frac{\partial V}{\partial x} \right\rangle
    \end{align}
    这正是经典力学中 $\dot{x} = p/m$ 和 $\dot{p} = -\partial V/\partial x$ 的量子对应——\textbf{埃伦费斯特定理表明，力学量的期望值服从经典运动方程}。
\end{example}

\subsubsection{\protect\hyperlink{:}{维里定理}}
\addtocontents{toc}{\protect\hypertarget{:}{}}

维里定理给出了束缚态中平均动能与平均势能之间的关系。

\begin{theorem}[维里定理]
    若势场 $V(\vec{r})$ 是 $\vec{r}$ 的 $\nu$ 次齐次函数，即
    \begin{equation}
        V(\lambda\vec{r}) = \lambda^\nu V(\vec{r})
    \end{equation}
    则在束缚定态 $\ket{\Psi_n}$ 上，平均动能与平均势能满足
    \begin{equation}
        2\langle T \rangle = \nu\langle V \rangle
    \end{equation}
\end{theorem}

\begin{proof}
    考虑算符 $\hat{G} = \hat{x}\hat{p}$（一维情况），其期望值为 $\langle \hat{G} \rangle = \langle \hat{x}\hat{p} \rangle$。

    由埃伦费斯特定理，在定态中 $\frac{\mathrm{d}}{\mathrm{d}t}\langle \hat{G} \rangle = 0$，所以
    \begin{equation}
        0 = \frac{1}{i\hbar}\langle [\hat{x}\hat{p}, \hat{H}] \rangle
    \end{equation}
    计算对易子
    \begin{equation}
        [\hat{x}\hat{p}, \hat{H}] = [\hat{x}\hat{p}, \frac{\hat{p}^2}{2m}] + [\hat{x}\hat{p}, V(\hat{x})]
    \end{equation}
    利用 $[\hat{x}, \hat{p}] = i\hbar$
    \begin{align}
        [\hat{x}\hat{p}, \hat{p}^2] &= \hat{x}[\hat{p}, \hat{p}^2] + [\hat{x}, \hat{p}^2]\hat{p} = 0 + 2i\hbar\hat{p}^2 \\
        [\hat{x}\hat{p}, V(\hat{x})] &= \hat{x}[\hat{p}, V] + [\hat{x}, V]\hat{p} = -i\hbar\hat{x}V' + 0
    \end{align}
    因此
    \begin{equation}
        0 = \frac{1}{i\hbar}\left(\frac{i\hbar}{m}\langle \hat{p}^2 \rangle - i\hbar\langle \hat{x}V' \rangle\right) = \frac{2}{m}\langle T \rangle - \langle \hat{x}V' \rangle
    \end{equation}
    即 $2\langle T \rangle = \langle \hat{x}V' \rangle$。

    由齐次函数的欧拉定理，$\nu V = \vec{r}\cdot\nabla V$，所以 $\langle \hat{x}V' \rangle = \nu\langle V \rangle$，代入即得
    \begin{equation}
        2\langle T \rangle = \nu\langle V \rangle
    \end{equation}
\end{proof}

\begin{remark}
    对于库仑势（$\nu = -1$），$2\langle T \rangle = -\langle V \rangle$，所以 $\langle T \rangle = -E$，$\langle V \rangle = 2E$（$E < 0$ 为束缚态能量）。对于谐振子势（$\nu = 2$），$\langle T \rangle = \langle V \rangle$——能量在动能和势能之间平均分配。
\end{remark}

\subsubsection{\protect\hyperlink{:}{Hellmann-Feynman定理}}
\addtocontents{toc}{\protect\hypertarget{:}{}}

\begin{theorem}[Hellmann-Feynman定理]
    设哈密顿量 $\hat{H}(\lambda)$ 含有参数 $\lambda$，束缚定态能量为 $E_n(\lambda)$，则
    \begin{equation}
        \frac{\partial E_n}{\partial \lambda} = \left\langle \frac{\partial \hat{H}}{\partial \lambda} \right\rangle_n
    \end{equation}
    其中右侧是在第 $n$ 个本征态上的期望值。
\end{theorem}

\begin{proof}
    由本征方程 $\hat{H}\ket{\psi_n} = E_n\ket{\psi_n}$，对 $\lambda$ 求导
    \begin{equation}
        \frac{\partial\hat{H}}{\partial\lambda}\ket{\psi_n} + \hat{H}\frac{\partial\ket{\psi_n}}{\partial\lambda} = \frac{\partial E_n}{\partial\lambda}\ket{\psi_n} + E_n\frac{\partial\ket{\psi_n}}{\partial\lambda}
    \end{equation}
    左乘 $\bra{\psi_n}$
    \begin{equation}
        \bra{\psi_n}\frac{\partial\hat{H}}{\partial\lambda}\ket{\psi_n} + \bra{\psi_n}\hat{H}\frac{\partial\ket{\psi_n}}{\partial\lambda} = \frac{\partial E_n}{\partial\lambda}\braket{\psi_n}{\psi_n} + E_n\bra{\psi_n}\frac{\partial\ket{\psi_n}}{\partial\lambda}
    \end{equation}
    由 $\bra{\psi_n}\hat{H} = E_n\bra{\psi_n}$（厄米性），后两项相消，得
    \begin{equation}
        \frac{\partial E_n}{\partial\lambda} = \bra{\psi_n}\frac{\partial\hat{H}}{\partial\lambda}\ket{\psi_n} = \left\langle \frac{\partial\hat{H}}{\partial\lambda} \right\rangle_n
    \end{equation}
\end{proof}

\begin{remark}
    H-F 定理在计算能级时非常有用。例如，对于氢原子，若将 $e^2$ 视为参数 $\lambda$，则
    \begin{equation}
        \frac{\partial E_n}{\partial(e^2)} = \left\langle \frac{\partial\hat{H}}{\partial(e^2)} \right\rangle_n = -\left\langle \frac{1}{4\pi\varepsilon_0 r} \right\rangle_n
    \end{equation}
    这给出了 $\langle 1/r \rangle$ 与能级的关系。
\end{remark}

\subsection{\protect\hyperlink{:}{常用的对易关系}}
\addtocontents{toc}{\protect\hypertarget{:}{}}

以下对易关系在量子力学中反复出现，需要熟练掌握。

\subsubsection{\protect\hyperlink{:}{基本对易关系}}
\addtocontents{toc}{\protect\hypertarget{:}{}}

\begin{theorem}[坐标-动量对易关系]
    \begin{equation}
        [\hat{x}_i, \hat{p}_j] = i\hbar\delta_{ij}
    \end{equation}
    其中 $i, j = 1, 2, 3$ 对应三个空间分量。
\end{theorem}

\begin{proof}
    在坐标表象中，$\hat{x}_i = x_i$，$\hat{p}_j = -i\hbar\frac{\partial}{\partial x_j}$。对任意波函数 $\psi(\vec{r})$
    \begin{equation}
        [\hat{x}_i, \hat{p}_j]\psi = x_i(-i\hbar\frac{\partial\psi}{\partial x_j}) - (-i\hbar\frac{\partial}{\partial x_j})(x_i\psi) = -i\hbar x_i\frac{\partial\psi}{\partial x_j} + i\hbar\left(\delta_{ij}\psi + x_i\frac{\partial\psi}{\partial x_j}\right) = i\hbar\delta_{ij}\psi
    \end{equation}
\end{proof}

\subsubsection{\protect\hyperlink{:}{角动量对易关系}}
\addtocontents{toc}{\protect\hypertarget{:}{}}

\begin{theorem}[角动量分量之间的对易关系]
    \begin{equation}
        [\hat{L}_i, \hat{L}_j] = i\hbar\varepsilon_{ijk}\hat{L}_k
    \end{equation}
    其中 $\varepsilon_{ijk}$ 是 Levi-Civita 符号。
\end{theorem}

\begin{proof}
    由 $\hat{L}_i = \varepsilon_{ijk}\hat{x}_j\hat{p}_k$，以 $[\hat{L}_x, \hat{L}_y]$ 为例
    \begin{align}
        [\hat{L}_x, \hat{L}_y] &= [\hat{y}\hat{p}_z - \hat{z}\hat{p}_y, \hat{z}\hat{p}_x - \hat{x}\hat{p}_z] \nonumber \\
        &= [\hat{y}\hat{p}_z, \hat{z}\hat{p}_x] - [\hat{y}\hat{p}_z, \hat{x}\hat{p}_z] - [\hat{z}\hat{p}_y, \hat{z}\hat{p}_x] + [\hat{z}\hat{p}_y, \hat{x}\hat{p}_z]
    \end{align}
    利用 $[\hat{x}_i, \hat{p}_j] = i\hbar\delta_{ij}$，只有含 $[\hat{p}_z, \hat{z}]$ 和 $[\hat{z}, \hat{p}_z]$ 的项不为零
    \begin{equation}
        [\hat{L}_x, \hat{L}_y] = \hat{y}[\hat{p}_z, \hat{z}]\hat{p}_x + \hat{x}[\hat{z}, \hat{p}_z]\hat{p}_y = \hat{y}(-i\hbar)\hat{p}_x + \hat{x}(i\hbar)\hat{p}_y = i\hbar(\hat{x}\hat{p}_y - \hat{y}\hat{p}_x) = i\hbar\hat{L}_z
    \end{equation}
    其他分量类似。
\end{proof}

\begin{remark}
    由此可得 $[\hat{L}^2, \hat{L}_i] = 0$（对 $i = x, y, z$ 均成立）。这意味着 $\hat{L}^2$ 与任意一个角动量分量（通常取 $\hat{L}_z$）可以同时确定——这就是为什么氢原子用 $l, m$ 两个量子数来描述角动量。
\end{remark}

\subsubsection{\protect\hyperlink{:}{函数形式的对易关系}}
\addtocontents{toc}{\protect\hypertarget{:}{}}

\begin{theorem}[与函数的对易关系]
    \begin{align}
        [\hat{p}_x, f(\hat{x})] &= -i\hbar\frac{\partial f}{\partial x} \\
        [\hat{x}, g(\hat{p}_x)] &= i\hbar\frac{\partial g}{\partial p_x}
    \end{align}
\end{theorem}

\begin{proof}
    以第一个为例。在坐标表象中
    \begin{equation}
        [\hat{p}_x, f(x)]\psi = -i\hbar\frac{\partial(f\psi)}{\partial x} - f(-i\hbar\frac{\partial\psi}{\partial x}) = -i\hbar\frac{\partial f}{\partial x}\psi
    \end{equation}
    故 $[\hat{p}_x, f(x)] = -i\hbar f'(x)$。第二个同理，在动量表象中证明。
\end{proof}

\begin{remark}
    这两个公式的记忆口诀：$[\hat{p}, f] = -i\hbar f'$（"p 左乘 f 等于求导"），$[\hat{x}, g] = i\hbar g'$（类似）。注意 $i\hbar$ 的正负号。
\end{remark}







\section{\protect\hyperlink{:}{表象}}
\addtocontents{toc}{\protect\hypertarget{:}{}}

量子力学中，态矢量 $\ket{\psi}$ 生活在抽象的希尔伯特空间中。为了将其与可观测量联系起来，我们需要选择一组基底来"展开"态矢量——这就是\textbf{表象}的核心思想。

\subsection{\protect\hyperlink{:}{表象的基本概念}}
\addtocontents{toc}{\protect\hypertarget{:}{}}

\subsubsection{\protect\hyperlink{:}{什么是表象}}
\addtocontents{toc}{\protect\hypertarget{:}{}}

\begin{definition}[表象]
    给定一个厄米算符 $\hat{A}$ 及其正交归一的本征态系 $\{\ket{a_n}\}$，以这组本征态作为希尔伯特空间的基底来展开态矢量和力学量，称为\textbf{$\hat{A}$ 表象}。
\end{definition}

表象的本质就是选择一组完备基底。不同的表象对应不同的基底选择，就像同一个矢量可以在不同的坐标系中用不同的分量来表示。

\subsubsection{\protect\hyperlink{:}{态矢量在表象中的表示}}
\addtocontents{toc}{\protect\hypertarget{:}{}}

在 $\hat{A}$ 表象中，任意态矢量 $\ket{\psi}$ 可以展开为
\begin{equation}
    \ket{\psi} = \sum_n c_n\ket{a_n}
\end{equation}

其中展开系数为
\begin{equation}
    c_n = \braket{a_n}{\psi}
\end{equation}

这就是态矢量在 $\hat{A}$ 表象中的表示——一个由所有展开系数构成的列向量
\begin{equation}
    \ket{\psi} \longleftrightarrow
    \begin{pmatrix}
        c_1 \\ c_2 \\ \vdots \\ c_n
    \end{pmatrix}
    =
    \begin{pmatrix}
        \braket{a_1}{\psi} \\ \braket{a_2}{\psi} \\ \vdots \\ \braket{a_n}{\psi}
    \end{pmatrix}
\end{equation}

这里用到了\textbf{完备性关系}（封闭性关系）
\begin{equation}
    \sum_n \ket{a_n}\bra{a_n} = \hat{1}
\end{equation}

这个关系是表象理论的数学基础，它保证了任意态矢量都可以用本征态系来展开。

\subsubsection{\protect\hyperlink{:}{算符在表象中的表示}}
\addtocontents{toc}{\protect\hypertarget{:}{}}

算符也可以在表象中用矩阵表示。设 $\hat{B}$ 是一个算符，其在 $\hat{A}$ 表象中的矩阵元定义为
\begin{equation}
    B_{mn} = \bra{a_m}\hat{B}\ket{a_n}
\end{equation}

算符 $\hat{B}$ 作用在态矢量 $\ket{\psi}$ 上得到新态 $\ket{\phi} = \hat{B}\ket{\psi}$，展开后
\begin{equation}
    \phi_m = \bra{a_m}\hat{B}\ket{\psi} = \sum_n \bra{a_m}\hat{B}\ket{a_n}\braket{a_n}{\psi} = \sum_n B_{mn}c_n
\end{equation}

这正是矩阵乘法的形式
\begin{equation}
    \begin{pmatrix}
        \phi_1 \\ \phi_2 \\ \vdots
    \end{pmatrix}
    =
    \begin{pmatrix}
        B_{11} & B_{12} & \cdots \\
        B_{21} & B_{22} & \cdots \\
        \vdots & \vdots & \ddots
    \end{pmatrix}
    \begin{pmatrix}
        c_1 \\ c_2 \\ \vdots
    \end{pmatrix}
\end{equation}

\begin{remark}
    厄米算符在自身表象中的矩阵是对角矩阵，对角元就是本征值
    \begin{equation}
        A_{mn} = \bra{a_m}\hat{A}\ket{a_n} = a_n\braket{a_m}{a_n} = a_n\delta_{mn}
    \end{equation}
    这就是为什么求解本征值问题如此重要——它相当于找到一个表象使算符对角化。
\end{remark}

\subsection{\protect\hyperlink{:}{坐标表象与动量表象}}
\addtocontents{toc}{\protect\hypertarget{:}{}}

最常用的两种表象是坐标表象和动量表象。它们的本征值是连续的，因此展开用积分代替求和。

\subsubsection{\protect\hyperlink{:}{坐标表象}}
\addtocontents{toc}{\protect\hypertarget{:}{}}

\begin{definition}[坐标表象]
    以坐标算符 $\hat{x}$ 的本征态 $\ket{x}$ 为基底的表象。本征方程为
    \begin{equation}
        \hat{x}\ket{x} = x\ket{x}
    \end{equation}
    其中 $x$ 可取任意实数（连续谱），正交归一关系为
    \begin{equation}
        \braket{x}{x'} = \delta(x - x')
    \end{equation}
    完备性关系为
    \begin{equation}
        \int_{-\infty}^{+\infty}\mathrm{d}x\,\ket{x}\bra{x} = \hat{1}
    \end{equation}
\end{definition}

态矢量在坐标表象中的表示就是\textbf{波函数}
\begin{equation}
    \psi(x) = \braket{x}{\psi}
\end{equation}

波函数的物理意义由玻恩诠释给出：$|\psi(x)|^2$ 是粒子出现在 $x$ 处的概率密度。

\subsubsection{\protect\hyperlink{:}{动量表象}}
\addtocontents{toc}{\protect\hypertarget{:}{}}

\begin{definition}[动量表象]
    以动量算符 $\hat{p}$ 的本征态 $\ket{p}$ 为基底的表象。本征方程为
    \begin{equation}
        \hat{p}\ket{p} = p\ket{p}
    \end{equation}
    在坐标表象中，动量本征态为平面波
    \begin{equation}
        \braket{x}{p} = \frac{1}{\sqrt{2\pi\hbar}}e^{ipx/\hbar}
    \end{equation}
\end{definition}

态矢量在动量表象中的表示为
\begin{equation}
    \varphi(p) = \braket{p}{\psi}
\end{equation}

\subsubsection{\protect\hyperlink{:}{坐标表象与动量表象的关系}}
\addtocontents{toc}{\protect\hypertarget{:}{}}

坐标表象和动量表象通过\textbf{傅里叶变换}相联系。

\begin{theorem}[坐标-动量表象的傅里叶变换关系]
    一维情况下
    \begin{align}
        \psi(x) &= \frac{1}{\sqrt{2\pi\hbar}}\int_{-\infty}^{+\infty}\varphi(p)e^{ipx/\hbar}\mathrm{d}p \\
        \varphi(p) &= \frac{1}{\sqrt{2\pi\hbar}}\int_{-\infty}^{+\infty}\psi(x)e^{-ipx/\hbar}\mathrm{d}x
    \end{align}
\end{theorem}

\begin{proof}
    利用完备性关系
    \begin{equation}
        \psi(x) = \braket{x}{\psi} = \int\mathrm{d}p\,\braket{x}{p}\braket{p}{\psi} = \int\mathrm{d}p\,\frac{e^{ipx/\hbar}}{\sqrt{2\pi\hbar}}\varphi(p)
    \end{equation}
    逆变换同理可得。
\end{proof}

\begin{remark}
    坐标表象和动量表象是完全等价的——它们只是同一物理实在的两种不同描述方式。波函数 $\psi(x)$ 和 $\varphi(p)$ 包含相同的信息，只是编码方式不同。这正体现了量子力学中表象的对称性。
\end{remark}

\subsubsection{\protect\hyperlink{:}{不同表象中的薛定谔方程}}
\addtocontents{toc}{\protect\hypertarget{:}{}}

薛定谔方程在不同表象中有不同的形式。

在坐标表象中
\begin{equation}
    -\frac{\hbar^2}{2m}\frac{\mathrm{d}^2\psi}{\mathrm{d}x^2} + V(x)\psi = E\psi
\end{equation}

在动量表象中，对坐标表象的方程做傅里叶变换，得到
\begin{equation}
    \frac{p^2}{2m}\varphi(p) + \int_{-\infty}^{+\infty}V_{pp'}\varphi(p')\mathrm{d}p' = E\varphi(p)
\end{equation}

其中势能矩阵元为
\begin{equation}
    V_{pp'} = \frac{1}{2\pi\hbar}\int_{-\infty}^{+\infty}V(x)e^{i(p-p')x/\hbar}\mathrm{d}x
\end{equation}

\begin{remark}
    动量表象中的薛定谔方程不再是微分方程，而是积分方程。当 $V(x)$ 是 $x$ 的多项式时，$V_{pp'}$ 涉及对 $p$ 的导数（利用 $[\hat{x}, f(\hat{p})] = i\hbar\frac{\partial f}{\partial p}$），方程变为微分方程。例如，谐振子势 $V = \frac{1}{2}m\omega^2 x^2$ 在动量表象中变为
    \begin{equation}
        \frac{p^2}{2m}\varphi(p) - \frac{m\omega^2\hbar^2}{2}\frac{\mathrm{d}^2\varphi}{\mathrm{d}p^2} = E\varphi(p)
    \end{equation}
    与坐标表象形式完全对称（$x \leftrightarrow p/m\omega$）。
\end{remark}

\subsection{\protect\hyperlink{:}{表象变换}}
\addtocontents{toc}{\protect\hypertarget{:}{}}

同一个量子态在不同表象中有不同的表示，它们之间通过\textbf{表象变换}相联系。

\subsubsection{\protect\hyperlink{:}{变换矩阵}}
\addtocontents{toc}{\protect\hypertarget{:}{}}

设 $\hat{A}$ 表象的基底为 $\{\ket{a_i}\}$，$\hat{B}$ 表象的基底为 $\{\ket{b_\alpha}\}$。

\begin{definition}[变换矩阵]
    变换矩阵 $U$ 的矩阵元定义为
    \begin{equation}
        U_{\alpha i} = \braket{b_\alpha}{a_i}
    \end{equation}
    即 $U$ 的第 $(\alpha, i)$ 个元素是 $\hat{A}$ 表象的第 $i$ 个基矢在 $\hat{B}$ 表象中的第 $\alpha$ 个分量。
\end{definition}

利用完备性关系，两个表象的基底之间可以相互展开
\begin{align}
    \ket{a_i} &= \sum_\alpha \ket{b_\alpha}\braket{b_\alpha}{a_i} = \sum_\alpha U_{\alpha i}\ket{b_\alpha} \\
    \ket{b_\alpha} &= \sum_i \ket{a_i}\braket{a_i}{b_\alpha} = \sum_i (U^{-1})_{i\alpha}\ket{a_i}
\end{align}

其中 $(U^{-1})_{i\alpha} = \braket{a_i}{b_\alpha} = U_{\alpha i}^*$。

\subsubsection{\protect\hyperlink{:}{幺正性}}
\addtocontents{toc}{\protect\hypertarget{:}{}}

\begin{theorem}[变换矩阵的幺正性]
    变换矩阵 $U$ 是幺正矩阵，即
    \begin{equation}
        U^\dagger U = UU^\dagger = \hat{1}
    \end{equation}
    等价地，$U^{-1} = U^\dagger$。
\end{theorem}

\begin{proof}
    计算 $(U^\dagger U)_{ij}$
    \begin{equation}
        (U^\dagger U)_{ij} = \sum_\alpha (U^\dagger)_{i\alpha}U_{\alpha j} = \sum_\alpha U_{\alpha i}^* U_{\alpha j} = \sum_\alpha \braket{a_i}{b_\alpha}\braket{b_\alpha}{a_j} = \braket{a_i}{a_j} = \delta_{ij}
    \end{equation}
    故 $U^\dagger U = \hat{1}$。同理可证 $UU^\dagger = \hat{1}$。
\end{proof}

\begin{remark}
    幺正变换保持内积不变（即保持概率守恒）：若 $\ket{\psi'} = U\ket{\psi}$，则
    \begin{equation}
        \braket{\psi'}{\psi'} = \bra{\psi}U^\dagger U\ket{\psi} = \braket{\psi}{\psi}
    \end{equation}
\end{remark}

\subsubsection{\protect\hyperlink{:}{态矢量和算符的变换}}
\addtocontents{toc}{\protect\hypertarget{:}{}}

\begin{theorem}[态矢量的表象变换]
    若态矢量在 $\hat{A}$ 表象中的表示为 $\ket{\psi}_A$，在 $\hat{B}$ 表象中的表示为 $\ket{\psi}_B$，则
    \begin{equation}
        \ket{\psi}_B = U\ket{\psi}_A
    \end{equation}
\end{theorem}

\begin{proof}
    在 $\hat{B}$ 表象中，态矢量的第 $\alpha$ 个分量为
    \begin{equation}
        (\psi_B)_\alpha = \braket{b_\alpha}{\psi} = \sum_i \braket{b_\alpha}{a_i}\braket{a_i}{\psi} = \sum_i U_{\alpha i}(\psi_A)_i
    \end{equation}
    即 $\ket{\psi}_B = U\ket{\psi}_A$。
\end{proof}

\begin{theorem}[算符的表象变换]
    若算符 $\hat{F}$ 在 $\hat{A}$ 表象中的矩阵为 $F_A$，在 $\hat{B}$ 表象中的矩阵为 $F_B$，则
    \begin{equation}
        F_B = UF_AU^\dagger
    \end{equation}
\end{theorem}

\begin{proof}
    在 $\hat{B}$ 表象中，矩阵元为
    \begin{equation}
        (F_B)_{\alpha\beta} = \bra{b_\alpha}\hat{F}\ket{b_\beta} = \sum_{i,j}\braket{b_\alpha}{a_i}\bra{a_i}\hat{F}\ket{a_j}\braket{a_j}{b_\beta} = \sum_{i,j}U_{\alpha i}(F_A)_{ij}(U^\dagger)_{j\beta}
    \end{equation}
    即 $F_B = UF_AU^\dagger$。
\end{proof}

\begin{remark}
    幺正变换 $F_B = UF_AU^\dagger$ 保持算符的本征值不变。设 $\hat{F}\ket{\psi} = f\ket{\psi}$，则
    \begin{equation}
        F_B(U\ket{\psi}) = UF_AU^\dagger U\ket{\psi} = UF_A\ket{\psi} = f(U\ket{\psi})
    \end{equation}
    因此 $U\ket{\psi}$ 是 $F_B$ 的本征态，本征值同为 $f$。这说明\textbf{物理量的本征值与表象选择无关}——这是表象理论自洽性的关键。
\end{remark}

\subsection{\protect\hyperlink{:}{表象的选取原则}}
\addtocontents{toc}{\protect\hypertarget{:}{}}

在实际问题中，选择合适的表象可以大大简化计算。

\begin{enumerate}
    \item 若哈密顿量 $\hat{H}$ 有简明的本征态系（如无限深势阱、谐振子），选择能量表象（$\hat{H}$ 表象）最方便，此时 $\hat{H}$ 是对角矩阵。
    \item 若势能 $V(x)$ 形式简单，选择坐标表象最方便，薛定谔方程是微分方程。
    \item 若势能 $V(x)$ 是 $x$ 的多项式，选择动量表象有时更方便。
    \item 对于角动量问题，选择 $\{\hat{L}^2, \hat{L}_z\}$ 表象最自然。
    \item 对于自旋问题，选择 $\hat{S}_z$ 表象，态矢量是二维列向量。
\end{enumerate}

\begin{remark}
    表象的选择不会影响物理结果。不同表象之间的变换是幺正变换，所有可观测量的期望值、概率分布等物理量都保持不变。表象只是计算的工具，不是物理本身。
\end{remark}






\section{\protect\hyperlink{:}{三维定态问题}}
\addtocontents{toc}{\protect\hypertarget{:}{}}

一维问题帮助我们建立了量子力学的基本概念，但真实世界是三维的。三维问题中，角动量的量子化是一个核心主题。

\subsection{\protect\hyperlink{:}{中心力场}}
\addtocontents{toc}{\protect\hypertarget{:}{}}

当势能只依赖于到某一点的距离 $r$（即 $V = V(r)$）时，问题具有球对称性，称为\textbf{中心力场问题}。氢原子、类氢离子都是典型的中心力场问题。

\subsubsection{\protect\hyperlink{:}{分离变量}}
\addtocontents{toc}{\protect\hypertarget{:}{}}

在球坐标系 $(r, \theta, \varphi)$ 中，定态方程为
\begin{equation}
    -\frac{\hbar^2}{2m}\left[\frac{1}{r^2}\frac{\partial}{\partial r}\left(r^2\frac{\partial\psi}{\partial r}\right) + \frac{1}{r^2\sin\theta}\frac{\partial}{\partial\theta}\left(\sin\theta\frac{\partial\psi}{\partial\theta}\right) + \frac{1}{r^2\sin^2\theta}\frac{\partial^2\psi}{\partial\varphi^2}\right] + V(r)\psi = E\psi
\end{equation}

利用球对称性，可以分离变量
\begin{equation}
    \psi(r,\theta,\varphi) = R(r)Y_l^m(\theta,\varphi)
\end{equation}

其中 $Y_l^m(\theta,\varphi)$ 是\textbf{球谐函数}，满足角动量方程；$R(r)$ 满足径向方程。

\subsubsection{\protect\hyperlink{:}{角动量量子化}}
\addtocontents{toc}{\protect\hypertarget{:}{}}

角动量是中心力场问题中的守恒量。量子力学中，轨道角动量算符为
\begin{equation}
    \hat{L}^2 = -\hbar^2\left[\frac{1}{\sin\theta}\frac{\partial}{\partial\theta}\left(\sin\theta\frac{\partial}{\partial\theta}\right) + \frac{1}{\sin^2\theta}\frac{\partial^2}{\partial\varphi^2}\right]
\end{equation}

$\hat{L}^2$ 和 $\hat{L}_z$ 的共同本征函数就是球谐函数 $Y_l^m(\theta,\varphi)$，满足
\begin{align}
    \hat{L}^2 Y_l^m &= l(l+1)\hbar^2 Y_l^m \\
    \hat{L}_z Y_l^m &= m\hbar Y_l^m
\end{align}

其中 $l = 0, 1, 2, \ldots$ 为\textbf{角量子数}，$m = -l, -l+1, \ldots, l-1, l$ 为\textbf{磁量子数}。

角动量的大小和 $z$ 分量都是量子化的
\begin{equation}
    |\vec{L}| = \sqrt{l(l+1)}\hbar, \qquad L_z = m\hbar
\end{equation}

\subsubsection{\protect\hyperlink{:}{氢原子}}
\addtocontents{toc}{\protect\hypertarget{:}{}}

氢原子是最简单的原子系统，其势能为库仑势
\begin{equation}
    V(r) = -\frac{e^2}{4\pi\varepsilon_0 r}
\end{equation}

求解径向方程，可以得到能量本征值为
\begin{equation}
    \boxed{E_n = -\frac{13.6 \text{ eV}}{n^2}, \qquad n = 1, 2, 3, \ldots}
\end{equation}

其中 $n$ 为\textbf{主量子数}。对于给定的 $n$，角量子数的取值为 $l = 0, 1, 2, \ldots, n-1$。

完整的波函数为
\begin{equation}
    \psi_{nlm}(r,\theta,\varphi) = R_{nl}(r)Y_l^m(\theta,\varphi)
\end{equation}

三个量子数的物理意义：
\begin{enumerate}
    \item 主量子数 $n$：决定能量 $E_n$ 和电子轨道的大小。
    \item 角量子数 $l$：决定角动量的大小 $|\vec{L}| = \sqrt{l(l+1)}\hbar$ 和轨道的形状。习惯上用字母表示：$s(l=0)$、$p(l=1)$、$d(l=2)$、$f(l=3)$。
    \item 磁量子数 $m$：决定角动量的 $z$ 分量 $L_z = m\hbar$ 和轨道的空间取向。
\end{enumerate}

\begin{remark}
    氢原子的能级只依赖于主量子数 $n$，与 $l$、$m$ 无关——这是库仑势的特殊对称性（$SO(4)$ 对称性）的结果。对于更复杂的原子（多电子原子），能级会依赖于 $l$，出现\textbf{能级分裂}。
\end{remark}

\subsection{\protect\hyperlink{:}{自旋}}
\addtocontents{toc}{\protect\hypertarget{:}{}}

除了轨道角动量之外，电子还具有一种内禀的角动量——\textbf{自旋}。自旋是纯粹的量子力学概念，没有经典对应。

\subsubsection{\protect\hyperlink{:}{自旋角动量}}
\addtocontents{toc}{\protect\hypertarget{:}{}}

电子的自旋量子数为 $s = 1/2$，自旋角动量的大小为
\begin{equation}
    |\vec{S}| = \sqrt{s(s+1)}\hbar = \frac{\sqrt{3}}{2}\hbar
\end{equation}

自旋的 $z$ 分量只能取两个值
\begin{equation}
    S_z = m_s\hbar, \qquad m_s = \pm\frac{1}{2}
\end{equation}

分别对应"自旋向上"（$m_s = +1/2$）和"自旋向下"（$m_s = -1/2$）。

\subsubsection{\protect\hyperlink{:}{自旋的矩阵表示}}
\addtocontents{toc}{\protect\hypertarget{:}{}}

在 $S_z$ 表象中，自旋态用二维列向量表示
\begin{equation}
    \ket{\uparrow} = \begin{pmatrix} 1 \\ 0 \end{pmatrix}, \qquad \ket{\downarrow} = \begin{pmatrix} 0 \\ 1 \end{pmatrix}
\end{equation}

自旋算符用 $2\times 2$ 的泡利矩阵表示
\begin{equation}
    \hat{S}_x = \frac{\hbar}{2}\sigma_x, \qquad \hat{S}_y = \frac{\hbar}{2}\sigma_y, \qquad \hat{S}_z = \frac{\hbar}{2}\sigma_z
\end{equation}

其中泡利矩阵为
\begin{equation}
    \sigma_x = \begin{pmatrix} 0 & 1 \\ 1 & 0 \end{pmatrix}, \qquad
    \sigma_y = \begin{pmatrix} 0 & -i \\ i & 0 \end{pmatrix}, \qquad
    \sigma_z = \begin{pmatrix} 1 & 0 \\ 0 & -1 \end{pmatrix}
\end{equation}

泡利矩阵满足对易关系 $[\sigma_i, \sigma_j] = 2i\varepsilon_{ijk}\sigma_k$ 和反对易关系 $\{\sigma_i, \sigma_j\} = 2\delta_{ij}$。

\subsection{\protect\hyperlink{:}{全同粒子与泡利不相容原理}}
\addtocontents{toc}{\protect\hypertarget{:}{}}

\subsubsection{\protect\hyperlink{:}{全同粒子的不可分辨性}}
\addtocontents{toc}{\protect\hypertarget{:}{}}

在量子力学中，全同粒子（如所有电子）是完全不可分辨的——交换任意两个全同粒子，系统的物理性质不变。

对于由 $N$ 个全同粒子组成的系统，交换第 $i$ 个和第 $j$ 个粒子后，波函数要么不变（\textbf{对称}），要么变号（\textbf{反对称}）
\begin{equation}
    \psi(\ldots, \vec{r}_i, \ldots, \vec{r}_j, \ldots) = \pm\psi(\ldots, \vec{r}_j, \ldots, \vec{r}_i, \ldots)
\end{equation}

\subsubsection{\protect\hyperlink{:}{费米子与玻色子}}
\addtocontents{toc}{\protect\hypertarget{:}{}}

\begin{enumerate}
    \item \textbf{费米子}：自旋为半整数（$s = 1/2, 3/2, \ldots$）的粒子，如电子、质子、中子。多粒子波函数是反对称的。
    \item \textbf{玻色子}：自旋为整数（$s = 0, 1, 2, \ldots$）的粒子，如光子、$\pi$ 介子。多粒子波函数是对称的。
\end{enumerate}

\subsubsection{\protect\hyperlink{:}{泡利不相容原理}}
\addtocontents{toc}{\protect\hypertarget{:}{}}

\begin{theorem}[泡利不相容原理]
    在一个量子系统中，不能有两个费米子处于完全相同的量子态。即两个电子不能同时具有相同的 $n$、$l$、$m$、$m_s$。
\end{theorem}

泡利不相容原理是理解原子结构和元素周期表的关键。它解释了为什么电子会逐层填充原子壳层，而不是全部挤在最低能级上。

\begin{remark}
    泡利不相容原理的深层原因是费米子波函数的反对称性。如果两个费米子处于相同的量子态，交换它们后波函数应该变号，但同时波函数又不变（因为是同一个态）——唯一的可能是波函数为零，即这种状态不存在。
\end{remark}

